\chapter*{致\quad 谢}
\thispagestyle{thesis-main}
\addcontentsline{toc}{chapter}{\texorpdfstring{致\quad 谢}{致 谢}}
\markboth{成都信息工程大学学士学位论文}{成都信息工程大学学士学位论文}

\thesisbodyfont
\justifying

在本课题完成过程中，我得到了老师、同学和家人的关心与帮助。在此，我向所有给予指导、支持和鼓励的人表示诚挚感谢。

衷心感谢指导教师何长春讲师。何老师在选题确定、研究思路梳理、系统实现和论文撰写过程中给予了耐心指导。也感谢人工智能学院各位老师在课程学习、科研训练和毕业设计阶段提供的帮助，使我能够将生成式人工智能、医学图像处理和平台开发等知识应用到本课题中。

感谢数派（成都）数据科技有限公司在毕业设计过程中提供的算力支持。StyleGAN2 和 3DGS 的训练对显存、计算时间和运行环境稳定性都有较高要求。相关算力资源为模型训练、结果生成和实验复核提供了重要保障，也使本文能够较完整地完成从二维 MRI 切片生成到三维表达与平台展示的实验流程。

感谢家人和同学在学习与生活中的陪伴、鼓励与支持。正是这些帮助让我能够在毕业设计过程中持续推进系统开发、实验整理和论文写作，并最终完成本课题研究。\par

\vspace{0.5\baselineskip}

\noindent 作者简介：\par

\begingroup
\renewcommand{\arraystretch}{1.4}
\noindent\begin{tabular*}{\textwidth}{@{}p{0.46\textwidth}@{\extracolsep{\fill}}p{0.46\textwidth}@{}}
姓名：黎曜玮 & 性别：男 \\
出生年月：{\tnr 2004}年{\tnr 7}月{\tnr 24}日 & 民族：汉族 \\
\multicolumn{2}{@{}l@{}}{E-mail：{\tnr 17781088366@163.com}} \\
\end{tabular*}
\endgroup\par
